\documentclass[12pt]{article}
\usepackage[margin=1in]{geometry}
\usepackage{enumitem}
\usepackage{titlesec}
\usepackage{parskip}
\usepackage{graphicx}
\usepackage{xcolor}
\usepackage{hyperref}
\usepackage{fontenc}

\title{\textbf{Bartels and Kramon 2020 RG Notes}\\
\large Bartels, Brandon L., and Eric Kramon. ``Does Public Support for Judicial Power Depend on Who is in Political Power? Testing a Theory of Partisan Alignment in Africa'' [2020].\\
\textit{American Political Science Review} 114(1): 144-163.}
\author{}
\date{}

\begin{document}

\maketitle

differentiation between judicial power and institutional legitimacy

``It is important to distinguish judicial power (and related concepts such as trust) from legitimacy because conflating them can lead to analyses that suffer from an observational equivalence problem: People may have high support for judicial power because (1) they agree with a court's rulings or expect partisan advantage or (2) they believe that a court has rightful authority regardless of partisan considerations. But those levels of support are derived from two drastically different processes. The first connotes instrumental support, while the latter connotes legitimacy (see also Bartels and Johnston 2020).'' (Bartels and Kramon, 2020, p. 146)

``Vertical power is judicial power over the mass public: When courts rule, the people should comply with those rulings. Horizontal power is judicial power over incumbent politicians: When the court makes a ruling that reverses government policy, the government should comply with that ruling.3'' (Bartels and Kramon, 2020, p. 146)

\end{document}
